\paragraph{From ternary coding to the marked flag.}
\label{inc-ampl-bandera}
The regions also provide the codes $w_\pi,w_e,w_\varphi$ and
$w_{\rm APP}$ in the inherited frame. The map
\[
 \gamma_{A_W}:w\longmapsto(w,wA_W)
\]
preserves the six-component word as its first coordinate and adds
its transformation by $A_W$. The relation $A_W^2=-I_6$ fixes
the quadratic compatibility of this realization. The next step,
$c\mapsto\operatorname{supp}(c)$, preserves the nonzero positions.
For weight-six words, the support identifies exactly the pair
$\{c,-c\}$, as proved by the minimum distance of the code. In this way,
the quotient by sign produces the hexads, while the complete code
remains available for operations requiring its symbols.

In this chart, the two constructions of the tetrad coincide:
\begin{equation}
 \{i:K_i\geq729\}
 =H_e\cap P_\pi=B_K=\{4,6,9,10\}.
 \label{inc-ampl-cuaterna}
\end{equation}
The left-hand side arises from an integer comparison of the components
of $K$; the right-hand side, from the intersection of the coded supports of
propagation and closure. The threshold belongs to the declared block
chart and remains explicit in the reader's definition. The
complete record preserves the values on which it acts, even when
different thresholds produce the same subset. The equality therefore
expresses compatibility between specified maps.

The support $H^-_\alpha$ completes that tetrad to a hexad.
The membership condition in support $P_\pi$, together with intersection
sizes $(4,3,2)$ relative to $H_e,H_\varphi,H_{\rm APP}$,
selects a unique hexad. Proposition~\ref{exc:selector-bandera}
proves that determination by the sign of $Z_0$ agrees with determination
by incidence; Lemma~\ref{exc:origen} fixes the origin in the resulting
frame. This information is organized as
\begin{equation}
 \widetilde{\mathcal I}_*
 =\bigl(B_K\subset H^-_\alpha;
         P_\pi,H_e,H_\varphi,H_{\rm APP}\bigr).
 \label{inc-ampl-marco}
\end{equation}
The inclusion is the flag; the four additional supports constitute
its marked frame. Simultaneous permutations transport the frame and
preserve its inclusions and intersections. The relevant information resides
in this transported configuration and its isomorphism class.

The five regions of $\pi$ here recover their internal differentiation.
They share $w_\pi$ as visible ternary code, but retain distinct emissions,
signatures, and multiplicities. In the binary realization on a
marked dodecad, the transversal region corresponds to the octad whose
intersection with the dodecad is the tetrad; the other four regions
correspond to lifts of the four hexads of its star. The
negative region is distinguished by $H^-_\alpha$, and chart order
individualizes the three positive regions. This correspondence preserves
regional roles and the distribution $432+4\cdot144=1008$, in addition
to realizing the incidence partition $1+4$.
